\documentclass[11pt]{article}
\usepackage[margin=1.6cm]{geometry}
\usepackage{amsmath}
\usepackage{tikz}
\usetikzlibrary{arrows.meta, positioning, shapes.geometric, calc, graphs, quotes}
\usepackage[most]{tcolorbox}
\tcbset{listing side text, size=small, fonttitle=\bfseries,
        colback=black!2, colframe=black!40, listing options={
        basicstyle=\ttfamily\footnotesize, breaklines=true}}

% ================= SHARED BASE (all approaches) =================
\makeatletter
\tikzset{
  tn/.style={inner sep=1pt, baseline=-.25em},
  copy/.style={circle, draw, inner sep=0pt, outer sep=0pt, minimum size=3pt},
  flat right/.style={regular polygon, regular polygon sides=3,
    shape border rotate=270, scale=.5, inner sep=3pt, draw},
  flat left/.style={regular polygon, regular polygon sides=3,
    shape border rotate=90, scale=.5, inner sep=3pt, draw},
  dmark/.style={{Circle[length=2.6pt]}-},
  every leg/.style={},
  leg length/.initial=.35,
}
% ---------- Approach A: legs as node options ----------
% legs={<angle>, <angle>/<label>, <angle>*, ...}   (* = derivative dot)
\tikzset{
  legs/.style={append after command={\pgfextra{\tn@legs{#1}}}},
}
\def\tn@legs#1{%
  \global\tn@dseenfalse
  \foreach \tn@item in {#1}{\expandafter\tn@leg\tn@item//\@nil}%
  \iftn@dseen \fill (\tikzlastnode.north east) circle (1.2pt);\fi}
\def\tn@leg#1/#2/#3\@nil{%
  \tn@checkstar#1*\@nil
  \iftn@d % derivative whisker from the SINGLE dot; bends toward its side
    \global\tn@dseentrue
    \ifdim\tn@ang pt<90pt
      \draw[every leg] (\tikzlastnode.north east) .. controls +(45:.12) .. ++(.28,.18);%
    \else
      \draw[every leg] (\tikzlastnode.north east) .. controls +(80:.12) .. ++(-.2,.28);%
    \fi
  \else
    \if\relax\detokenize{#2}\relax
      \draw[every leg] (\tikzlastnode) -- ++(\tn@ang:\pgfkeysvalueof{/tikz/leg length});%
    \else
      \pgfmathsetmacro\tn@anch{\tn@ang+180}%
      \draw[every leg] (\tikzlastnode) -- ++(\tn@ang:\pgfkeysvalueof{/tikz/leg length})
         node[font=\tiny, anchor=\tn@anch, inner sep=1.5pt] {$#2$};%
    \fi
  \fi}
\newif\iftn@d
\newif\iftn@dseen
\def\tn@checkstar#1*#2\@nil{%
  \def\tn@ang{#1}%
  \if\relax\detokenize{#2}\relax \tn@dfalse\else \tn@dtrue\fi}
\makeatother

% ---------- Approach C: mini-DSL ----------
% \tn{-a-{X'}-b-}  : dash-separated chain; leading/trailing dash = free end;
%  * = copy dot, > = right triangle, < = left triangle, trailing ' = d-legs
\ExplSyntaxOn
\NewDocumentCommand{\tn}{ O{2.0} m }
 {
  \begin{tikzpicture}[tn]
  \seq_set_split:Nnn \l_tmpa_seq { - } { #2 }
  \int_zero:N \l_tmpa_int
  \seq_map_inline:Nn \l_tmpa_seq
   {
    \int_incr:N \l_tmpa_int
    \tl_set:Nx \l_tmpa_tl { ( \fp_eval:n { (\l_tmpa_int - 1) * #1 } em , 0 ) }
    \str_case:nnF { ##1 }
     {
      { }  { \tn_place:nn { coordinate } { } }
      { * }{ \tn_place:nn { copy } { } }
      { > }{ \tn_place:nn { flat~right } { } }
      { < }{ \tn_place:nn { flat~left } { } }
     }
     { \tn_letter:n { ##1 } }
    \int_compare:nNnT { \l_tmpa_int } > { 1 }
     { \draw (n\int_eval:n{\l_tmpa_int-1}) -- (n\int_use:N\l_tmpa_int); }
   }
  \end{tikzpicture}
 }
\cs_new_protected:Npn \tn_place:nn #1 #2
 {
  \use:x { \exp_not:N \node [ #1 ] ( n \int_use:N \l_tmpa_int )
           at \tl_use:N \l_tmpa_tl } { #2 } ;
 }
\cs_new_protected:Npn \tn_letter:n #1
 {
  \tl_set:Nn \l_tmpb_tl { #1 }
  \tl_if_in:NnTF \l_tmpb_tl { ' }
   {
    \tl_remove_all:Nn \l_tmpb_tl { ' }
    \use:x { \exp_not:N \node [ legs = { 135 * , 45 * } ]
             ( n \int_use:N \l_tmpa_int ) at \tl_use:N \l_tmpa_tl }
      { $ ( \tl_use:N \l_tmpb_tl ) $ } ;
   }
   {
    \use:x { \exp_not:N \node ( n \int_use:N \l_tmpa_int )
             at \tl_use:N \l_tmpa_tl } { $ #1 $ } ;
   }
 }
\ExplSyntaxOff

% ---------- Approach D: pics (molecules) ----------
\tikzset{
  pics/kron/.style 2 args={code={
    \node[flat left]  (-in)  at (0,0) {};
    \node (-A) at (.55,.22) {$#1$};
    \node (-B) at (.55,-.22) {$#2$};
    \node[flat right] (-out) at (1.1,0) {};
    \draw (-in) -- (-A) -- (-out);
    \draw (-in) -- (-B) -- (-out);
  }},
  pics/det stack/.style={code={
    \node (-a) at (0,0) {$#1$};
    \node (-dots) at (.55,0) {$\cdots$};
    \node (-b) at (1.1,0) {$#1$};
    \draw (-a.north) -- ++(0,.18) coordinate (-at);
    \draw (-b.north) -- ++(0,.18) coordinate (-bt);
    \draw (-a.south) -- ++(0,-.18) coordinate (-ab);
    \draw (-b.south) -- ++(0,-.18) coordinate (-bb);
    \draw[line width=1.6pt] (-at -| -a.west) -- (-bt -| -b.east);
    \draw[line width=1.6pt] (-ab -| -a.west) -- (-bb -| -b.east);
  }},
  pics/tt chain/.style={code={
    \foreach \g [count=\c] in {#1} {
      \node (-g\c) at (\c*1.1, 0) {$\g$};
      \draw (-g\c) -- ++(0,-.35) node[font=\tiny, below, inner sep=1pt] {$i_\c$};
      \ifnum\c>1 \draw (-g\the\numexpr\c-1) -- (-g\c); \fi
    }
  }},
}

\pagestyle{empty}
\setlength\parindent{0pt}
\begin{document}

\section*{Four approaches to a tensor-diagram library}
Same test diagrams in each approach. Shared base: \texttt{tn} picture style,
\texttt{copy}, \texttt{flat left/right} node styles.

% =====================================================================
\subsection*{A. Legs as node options \hfill \normalfont\small (plain tikz + two new keys)}

\begin{tcblisting}{title=A1: derivative of $a^T X b$}
\begin{tikzpicture}[tn]
  \node[legs={180}] (a) {$a$};
  \node[right=.8em of a, legs={135*,45*}] (X) {$(X)$};
  \node[right=.8em of X, legs={0}]   (b) {$b$};
  \draw (a) -- (X) -- (b);
\end{tikzpicture}
\end{tcblisting}

\begin{tcblisting}{title=A2: TT chain with labeled legs}
\begin{tikzpicture}[tn]
  \node[legs={270/i_1}] (G1) {$G_1$};
  \node[right=1em of G1, legs={270/i_2}] (G2) {$G_2$};
  \node[right=1em of G2, legs={270/i_3}] (G3) {$G_3$};
  \draw (G1) -- (G2) -- (G3);
\end{tikzpicture}
\end{tcblisting}

\begin{tcblisting}{title=A3: Khatri--Rao $A \ast B$}
\begin{tikzpicture}[tn]
  \node[flat left, legs={180}] (f) {};
  \node (A) at (.55,.22)  {$A$};
  \node (B) at (.55,-.22) {$B$};
  \node[copy, legs={0}] (c) at (1.1,0) {};
  \draw (f) -- (A) -- (c);
  \draw (f) -- (B) -- (c);
\end{tikzpicture}
\end{tcblisting}

% =====================================================================
\subsection*{B. The \texttt{graphs} library \hfill \normalfont\small (chains as one expression)}

\begin{tcblisting}{title=B1: derivative of $a^T X b$}
\begin{tikzpicture}[tn]
\graph[grow right sep=.8em] {
  l[as=, coordinate] -- a[as=$a$] --
  X[as=$(X)$, legs={135*,45*}] --
  b[as=$b$] -- r[as=, coordinate]
};
\end{tikzpicture}
\end{tcblisting}

\begin{tcblisting}{title=B2: TT chain}
\begin{tikzpicture}[tn]
\graph[grow right sep=1em] {
  G1[as=$G_1$, legs={270/i_1}] --
  G2[as=$G_2$, legs={270/i_2}] --
  G3[as=$G_3$, legs={270/i_3}]
};
\end{tikzpicture}
\end{tcblisting}

\begin{tcblisting}{title=B3: Khatri--Rao (fan-out/fan-in)}
\begin{tikzpicture}[tn]
\graph[grow right sep=.6em,
       branch down sep=.5em] {
  l[as=, coordinate] --
  f[as=, flat left] --
  { A[as=$A$], B[as=$B$] } --
  c[as=, copy] -- r[as=, coordinate]
};
\end{tikzpicture}
\end{tcblisting}

% =====================================================================
\subsection*{C. A mini-DSL \hfill \normalfont\small (dashes are edges, like the book's inline notation)}

\begin{tcblisting}{title=C1: derivative of $a^T X b$}
\tn{-a-{X'}-b-}
\end{tcblisting}

\begin{tcblisting}{title=C2: matrix chain with copy dot}
\tn{-A-*-B-}
\end{tcblisting}

\begin{tcblisting}{title=C3: Khatri--Rao-ish chain}
\tn{-<-M-*-}
\end{tcblisting}

% =====================================================================
\subsection*{D. Pics: reusable molecules \hfill \normalfont\small (with anchors, so they compose)}

\begin{tcblisting}{title=D1: Kronecker bundle}
\begin{tikzpicture}[tn]
  \pic (k) {kron={A}{B}};
  \draw[double] (k-in) -- ++(-.35,0);
  \draw[double] (k-out) -- ++(.35,0);
\end{tikzpicture}
\end{tcblisting}

\begin{tcblisting}{title=D2: mixed product (511) by composition}
\begin{tikzpicture}[tn]
  \pic (k1) {kron={A}{B}};
  \pic (k2) at (1.6,0) {kron={C}{D}};
  \draw[double] (k1-in) -- ++(-.35,0);
  \draw[double] (k1-out) -- (k2-in);
  \draw[double] (k2-out) -- ++(.35,0);
\end{tikzpicture}
\end{tcblisting}

\begin{tcblisting}{title=D3: determinant stack + TT chain}
\begin{tikzpicture}[tn]
  \pic {det stack=X};
\end{tikzpicture}
\qquad
\begin{tikzpicture}[tn]
  \pic {tt chain={G_1,G_2,G_3}};
\end{tikzpicture}
\end{tcblisting}

\end{document}
